% Page 22 — Alban en sa « légende dorée »
% À intégrer dans le document principal (XeLaTeX).
% Nécessite le package lettrine.

\begin{center}
    {\Large\itshape Alban en sa « légende dorée »}
\end{center}

\medskip

\lettrine[lines=2]{L}{es} enfants venus faire leurs dévotions à la chapelle de Saint Joseph,
racontaient aux bons moines qui les recevaient, et ce, avec joie et un brin
de malice, les bruyants ébats qu’ils prenaient jusque dans les ruines du
vieux « castel ». Ils n’y rencontraient plus « le très méchant Alain de
Nougaret » dont on parlait encore aux veillées d’hiver comme d’un terrible
croquemitaine. « Il faut prier pour les pécheurs », se contentait de répondre
avec humilité le religieux pénitent.

\medskip

\lettrine[lines=2]{L}{es} paysans eux-mêmes qui venaient recommander leurs bestiaux à la
protection de Saint Roch dont la chapelle faisait pendant à celle de
Saint Joseph, ne se doutaient pas de la personnalité de frère Alain.
Il leur arrivait même de parler avec une vindicative rudesse de leur ancien
maître. « Heureusement qu’on a grillé Alain de Nougaret comme un poulet
et vraisemblablement il doit encore griller en enfer, car il avait vendu son
âme au diable, à ce que l’on dit ». « A tout pécheur miséricorde » se
contentait de répondre le frère bénédictin. Bref, après des années et des
années de prière et de pénitence, frère Alain termina sa course terrestre
et exhala son dernier soupir sous le scapulaire de Saint Benoît. Alors
seulement fut révélé aux uns et aux autres l’identité de ce religieux dont
la fin si édifiante avait fait quelque peu oublier la mauvaise réputation
du seigneur de Nougaret, baron d’Alban. Comme vous le
